% Appendix A, from protocol/uart_register_protocol.md: the contract both halves of the course are
% written against, printed whole, with its three parts as sections so the chapters can point at
% them as the lectures do ("Part 2 of the spec").

\chapter{The UART Register Protocol}
\label{proto:ch}

{\large\itshape The contract at the centre of the course.\par}

The FPGA side (Chapters~1 to~5) implements this contract as a peripheral behind the provided SPI
transport; the MCU side (Chapters~6 to~9) implements it as a master; and \chapref{10} proves the
two agree. Both halves are written against this appendix, not against each other's code. \textbf{If
an implementation and this appendix disagree, the implementation is wrong.}

There are two layers here, and they are independent:
\begin{enumerate}
  \item \textbf{The UART line protocol}: how the peripheral's \code{tx} and \code{rx} pins carry
    bytes to the outside world, the \emph{data plane}.
  \item \textbf{The register map and its SPI access}: how the ATmega328P reaches the peripheral's
    registers across the chip boundary, the \emph{control plane}.
\end{enumerate}

% ------------------------------------------------------------------------------------------
\section{Part 1: The UART Line Protocol}
\label{proto:part1}
This is the data plane: the peripheral's own asynchronous serial line, on \code{tx} and \code{rx}.

\begin{booktable}{@{}p{7.2em}L@{}}
  \thead{Parameter} & \thead{Value} \\
  \midrule
  Frame & 1 start bit (low), then the data bits, then an optional parity bit, then the stop bit
    or bits (high). \\
  Data bits & 8, \textbf{least significant first}. \\
  Idle & The line idles \textbf{high}. \\
  Parity & None, even or odd, selected by \code{CTRL}. Default none. \\
  Stop bits & 1 or 2, selected by \code{CTRL}. Default 1. \\
  Default frame & \textbf{8N1}: 8 data bits, no parity, 1 stop bit. \\
  Voltage & 3.3\,V logic, that of the DE0-CV. The data-plane peer must be 3.3\,V too, for
    example a CP2102 or an FT232 set to 3.3\,V. \\
\end{booktable}

\subsection{Baud rate}
\label{proto:baud}
The system clock is \textbf{50\,MHz}. The receiver oversamples the line at \textbf{16 times} the
baud rate, so the internal tick rate is \code{16 * baud}. \code{BAUD_DIV} holds the divider from
the 50\,MHz clock to that tick:

\begin{console}
BAUD_DIV = round( 50_000_000 / (16 * baud) )
\end{console}

\begin{narrowtable}{@{}rr@{}}
  \thead{Baud} & \thead{\code{BAUD_DIV}} \\
  \midrule
  9600 & 326 \\
  19200 & 163 \\
  57600 & 54 \\
  115200 & 27 \\
\end{narrowtable}

The transmitter emits one bit per 16 ticks; the receiver samples each bit at tick~8, mid-bit, after
detecting a start edge. \code{BAUD_DIV} is 16 bits, so with the oversampling factor of 16 the
lowest representable rate is \code{50_000_000 / (16 * 65535)}, about 48~baud.

Because \code{BAUD_DIV} is an integer, the achieved rate is only as close as the rounding allows.
At 115200 the exact divider is 27.126, so \code{BAUD_DIV = 27} gives 115740.7~baud, about 0.47\,\%
fast; at 9600 the error is $-0.15$\,\%. Both are comfortably inside what an 8-bit frame tolerates:
the receiver resynchronizes on every start bit, and only has to be still inside the stop bit when it
samples it, nine and a half bit periods later, which allows a few percent of combined error between
the two ends.

% ------------------------------------------------------------------------------------------
\section{Part 2: The Register Map}
\label{proto:part2}
This is the control plane: seven 32-bit registers, reached over the provided SPI transport of
Part~3. Only the low bits of each register are meaningful; the upper bits read back as zero and are
ignored on a write.

\begin{booktable}{@{}p{2.2em}p{7.2em}p{2.8em}p{2.4em}L@{}}
  \thead{Index} & \thead{Register} & \thead{Offset} & \thead{Access} & \thead{Description} \\
  \midrule
  0 & \code{STATUS} & \code{0x00} & R & Bit 0: TX ready (the TX FIFO is not full). Bit 1: RX valid
    (the RX FIFO is not empty). Bit 2: Error (any \code{ERROR_FLAGS} bit set). Bit 3: TX idle (the
    TX FIFO is empty and the line is idle). \\
  1 & \code{CTRL} & \code{0x04} & R/W & Bit 0: enable. Bits 2--1: parity (00 none, 01 even, 10 odd;
    11 reserved). Bit 3: stop bits (0 for 1, 1 for 2). Bit 4: RX-valid IRQ mask. Bit 5: TX-ready IRQ
    mask. \textbf{Reserved in this course: see \secref{proto:reserved}.} \\
  2 & \code{BAUD_DIV} & \code{0x08} & R/W & The baud divider (bits 15--0), per Part~1. \\
  3 & \code{TX_DATA} & \code{0x0C} & W & Bits 7--0: push one byte into the TX FIFO. \\
  4 & \code{RX_DATA} & \code{0x10} & R & Bits 7--0: the byte at the front of the RX FIFO. Reading
    it does \textbf{not} pop it. \\
  5 & \code{RX_POP} & \code{0x14} & W & Any write discards the front RX byte and advances the FIFO;
    the value is ignored. Write \code{0x1} by convention. \\
  6 & \code{ERROR_FLAGS} & \code{0x18} & R/W & Bit 0: framing error. Bit 1: parity error. Bit 2:
    overrun. Write \code{0x0} to clear all three. \\
\end{booktable}

Indices 7 to 15 are reserved: a read returns \code{0x00000000}, and a write is ignored.

\subsection{What this course implements, and what it reserves}
\label{proto:reserved}
The map above is the full contract. The peripheral built in Chapters~1 to~5 implements the part of
it that carries data, and stores but does not act on the rest:
\begin{itemize}
  \item \textbf{\code{CTRL} is stored and read back, but gates nothing.} The register bank has no
    \code{ctrl} output, so \code{CT_ENABLE}, the parity select, the stop-bit select and the two IRQ
    masks reach no logic. The peripheral is always enabled and always 8N1. The driver's
    \code{configure()} still writes the enable bit, so the two halves agree on the wire; the write
    simply has no effect on the hardware yet.
  \item \textbf{There is no interrupt line.} \code{uart_top} has eight ports and none of them is an
    IRQ output, so the two IRQ mask bits mask nothing. The system is poll-only, end to end.
  \item \textbf{Only the framing error flag has a producer.} The receiver emits \code{frame_err}
    and nothing else, so \code{ER_PARITY} and \code{ER_OVERRUN} are defined positions that read zero
    unless software writes a one to them. Parity detection and overrun detection are the natural
    extensions of \chapref{3} and \chapref{5} respectively.
\end{itemize}
These are reservations, not omissions: the positions are fixed so that an implementation which adds
them later stays compatible with both halves as written.

\subsection{The driver flow this map implies}
\label{proto:flow}
\begin{itemize}
  \item \textbf{Transmit}: poll \code{STATUS} bit 0 (TX ready); when it is set, write a byte to
    \code{TX_DATA}.
  \item \textbf{Receive}: poll \code{STATUS} bit 1 (RX valid); when it is set, read \code{RX_DATA},
    then write \code{RX_POP} to advance to the next byte.
\end{itemize}
This mirrors the poll-status, read-data, acknowledge shape of a typical hardware FIFO.

\subsection{Register semantics}
\label{proto:semantics}
This is what the register bank must implement. The TX and RX cores expose single-cycle pulses and
levels (Chapters~2 and~3); the register map promises poll-able bits, sticky error flags and
FIFO-backed data. Bridging the two is the register bank's whole job (\chapref{5}):
\begin{itemize}
  \item \textbf{\code{STATUS} bit 0 (TX ready)}: a level, meaning the TX FIFO is not full. A
    \code{TX_DATA} write while it is clear is dropped, because there is no room, which is why the
    driver must poll it first.
  \item \textbf{\code{STATUS} bit 1 (RX valid)}: a level, meaning the RX FIFO is not empty. It is
    set by the receiver pushing a byte, and cleared when \code{RX_POP} empties the FIFO.
  \item \textbf{\code{STATUS} bit 2 (Error) and \code{ERROR_FLAGS}}: the bank latches the receiver's
    framing, parity and overrun pulses into the three \code{ERROR_FLAGS} bits, and bit 2 of
    \code{STATUS} is their OR. An \code{ERROR_FLAGS} write of \code{0x0} clears the latches.
  \item \textbf{\code{TX_DATA}}: a write pushes exactly one byte into the TX FIFO. It is an
    \textbf{edge event}, one write for one byte, never a held level. On a transaction abort
    (Part~3), no byte is pushed.
  \item \textbf{\code{RX_DATA} and \code{RX_POP}}: \code{RX_DATA} is a \emph{pure read}. It returns
    the front byte and has no side effect, so it obeys the same latch-once rule as \code{STATUS}.
    Popping the FIFO is a \emph{separate write}, \code{RX_POP}, so that the side effect is a
    committed, abort-safe action and never a by-product of reading. This split is the point of
    \chapref{5}: \textbf{a read that popped a FIFO would need write-like commit-and-abort
    discipline; keeping the read pure and the pop explicit avoids that entirely.}
  \item \textbf{\code{BAUD_DIV} and \code{CTRL}}: plain read/write registers. \code{baud_gen} reads
    \code{BAUD_DIV} continuously, so a new divider takes effect from the next baud tick, in the
    middle of a frame if one is in flight; \code{CTRL} is stored only. Reconfiguring mid-frame is
    the driver's problem, not the bank's, which is what \code{STATUS} bit 3 (TX idle) is for: the
    driver waits for it before changing the baud rate.
\end{itemize}

% ------------------------------------------------------------------------------------------
\section{Part 3: The SPI Transport}
\label{proto:part3}
The register map above is reached over SPI by the provided \code{spi_slave} and
\code{spi_reg_bridge}. This section documents that transport fully: it is a black box in
\emph{this} course, which you use and do not write, but it is not magic.

\begin{booktable}{@{}p{7.2em}L@{}}
  \thead{Parameter} & \thead{Value} \\
  \midrule
  Roles & The MCU is the SPI master, the FPGA the slave. \\
  Mode & SPI mode 0 (\code{CPOL = 0}, \code{CPHA = 0}): \code{SCK} idles low, and both sides sample
    on the rising edge. \\
  Bit order & Most significant bit first, in every byte. \\
  \code{SCK} frequency & At most 1\,MHz. The Nano uses f\textsubscript{osc}/16, which is 1\,MHz. \\
  \code{SS} & Active low; one transaction per low period. \\
  Voltage & The Nano is 5\,V and the DE0-CV 3.3\,V, so every SPI line crosses a level shifter
    (\chapref{10}). \\
\end{booktable}

\subsection{Transactions}
\label{proto:txn}
Every transaction is exactly \textbf{5 bytes}: one command byte, then four data bytes. \code{SS}
falls before the command byte and rises after the fifth byte, staying low for the whole
transaction. The command byte is laid out as:

\begin{textdrawing}
Bit:      7    6    5    4    3    2    1    0
        +----+----+----+----+----+----+----+----+
        | W  | 0  | 0  | 0  |   register index  |
        +----+----+----+----+----+----+----+----+
\end{textdrawing}

\begin{itemize}
  \item \textbf{Bit 7 (\code{W})}: \code{1} for a write, \code{0} for a read.
  \item \textbf{Bits 6 to 4}: reserved, and must be \code{0}.
  \item \textbf{Bits 3 to 0}: the register index, which is the offset divided by 4, so 0 to 6 in the
    map above.
\end{itemize}
The two kinds of transaction, and the one way either can end early:
\begin{itemize}
  \item \textbf{Write (\code{W = 1})}: the master sends the 32-bit value in the four data bytes,
    most significant byte first. The slave commits it to the addressed register only when the fifth
    byte completes.
  \item \textbf{Read (\code{W = 0})}: at the end of the command byte the slave \textbf{latches the
    addressed register's value once}; the four data bytes clock it out on \code{MISO}, most
    significant byte first, while the master sends dummy \code{0x00} bytes. The latch-once rule is
    part of the contract: the four bytes are one coherent sample.
  \item \textbf{Abort}: if \code{SS} rises before the fifth byte, the transaction is abandoned with
    \textbf{no side effects}. Nothing commits, no \code{TX_DATA} push happens and no \code{RX_POP}.
    This is what makes the peripheral self-recovering.
\end{itemize}

% ------------------------------------------------------------------------------------------
\section{A Worked Example}
\label{proto:example}
Configuring 115200~8N1, sending \code{'H'} (\code{0x48}), then polling for and reading one received
byte. Each line is one transaction, as the master sends it; on a read, the four bytes after the
command byte are what comes back on \code{MISO}.

\begin{console}
Write BAUD_DIV : 82 00 00 00 1B   (0x1B = 27)
Write CTRL     : 81 00 00 00 01   (enable, 8N1)
Read  STATUS   : 00 xx xx xx xx   -> bit 0 set means TX ready
Write TX_DATA  : 83 00 00 00 48   ('H' pushed into the TX FIFO)
Read  STATUS   : 00 xx xx xx xx   -> bit 1 set means RX valid
Read  RX_DATA  : 04 00 00 00 XX   -> XX is the received byte
Write RX_POP   : 85 00 00 00 01   (advance the RX FIFO)
\end{console}
